% 导言区：按校赛示例页压缩标题间距，并保持正文为小四宋体、两字符缩进。

% 模板原标题使用宋体粗体，Windows 字体集会产生 SimSun 粗体回退警告。
% 按论文示例页收紧各级标题的段前、段后距离，并统一使用黑体。
% 一级标题采用四号，二级及以下标题采用正文小四字号。
\makeatletter
\renewcommand\section{\@startsection{section}{1}{\z@}%
  {-1.0ex \@plus -.2ex \@minus -.1ex}%
  {.8ex \@plus.1ex}%
  {\centering\normalfont\zihao{4}\heiti}}
\renewcommand\subsection{\@startsection{subsection}{2}{\z@}%
  {-1.0ex\@plus -.2ex \@minus -.1ex}%
  {.4ex \@plus .1ex}%
  {\normalfont\normalsize\heiti}}
\renewcommand\subsubsection{\@startsection{subsubsection}{3}{\z@}%
  {-1.0ex\@plus -.2ex \@minus -.1ex}%
  {.4ex \@plus .1ex}%
  {\normalfont\normalsize\heiti}}
\makeatother

\setlength{\parindent}{2em}
\setlength{\parskip}{0pt}
\renewcommand*{\baselinestretch}{1.38}
\graphicspath{{figures/}}

% 统一伪代码块：ruled 风格，支持中文与跨页。
\usepackage[most]{tcolorbox}
\newtcolorbox{pseudocodeblock}[1]{
  enhanced,
  colback=white,colframe=white,colbacktitle=white,coltitle=black,
  boxrule=0pt,arc=0pt,
  left=0pt,right=0pt,top=3pt,bottom=3pt,
  before skip=8pt,after skip=8pt,
  title={#1},fonttitle=\normalfont\normalsize,fontupper=\small,
  toptitle=2pt,bottomtitle=2pt,titlerule=.5pt,
  borderline north={.7pt}{0pt}{black},
  borderline south={.7pt}{0pt}{black}
}

\newcounter{pseudocodeline}
\newenvironment{pseudocodelines}{%
  \setcounter{pseudocodeline}{0}%
  \par\small\setlength{\parindent}{0pt}%
}{\par}
\newcommand{\pseudostep}[2][0em]{%
  \refstepcounter{pseudocodeline}%
  \noindent\makebox[2.2em][r]{\scriptsize\thepseudocodeline:}%
  \hspace{.8em}\hspace{#1}#2\par
}

% 图表题注直接使用黑体，避免宋体粗体回退。
\DeclareCaptionFont{hei}{\heiti}
\captionsetup[figure]{font={hei,minusfour},position=bottom}
\captionsetup[table]{font={hei,minusfour},position=top}

\newcommand{\vect}[1]{\boldsymbol{#1}}
